\def\level{Première}  % Classe à droite
\def\course{NSI}
\def\chaptername{Les Fonctions} % Titre au centre
\def\docname{AP02-S}        % Identifiant à gauche

\pagestyle{cours_FP}
\makeheader{} % Affiche le bandeau

\begin{easybox}{Gestion des entrées en Python}
    \begin{CodePythontexAlt}[Largeur=17cm,Centre,PremLigne=1,Lignes=false,]{}
        a = input("Entrer une valeur au clavier")           # a sera de type str
        b = int(input("Entrer une valeur au clavier"))      # b sera de type int
    \end{CodePythontexAlt}

    \begin{CodePythontexAlt}[Largeur=17cm,Centre,PremLigne=1,Lignes=false,]{}
         a = 6
         print(a)  # Affiche 6
         print("La valeur de a est", a)  # Affiche "La valeur de a est 6"   
      \end{CodePythontexAlt}
\end{easybox}	

\begin{methode}[Formatage de l'affichage avec \texttt{f-string} et \texttt{format()}]{}

    L'idée est de créer un texte à trou que l'on va compléter par des variables.
\begin{ConsolePythontex}[Label=false]{}
    name = 'Paul'
    age = 25
    print(f'Votre nom est {name} et vous avez {age} ans.')
\end{ConsolePythontex}  

\begin{ConsolePythontex}[Label=false]{}
    print('Votre nom est {} et vous avez {} ans'.format(name, age))
\end{ConsolePythontex}
\end{methode}


\begin{definition}[Les fonctions]{}

	\begin{CodePythontexAlt}[Centre,PremLigne=1]{}
		def f(x):
			""" Fonction renvoyant x ** 2 """
			y = x ** 2		# Calcul de x ** 2
			return y  		# valeur renvoyée: y
		\end{CodePythontexAlt}
		
   \end{definition}

		\begin{CodePythontexAlt}[Centre,PremLigne=1, Lignes=false,]{}
			a = 4   		# on affecte 4 à la variable a
			z = f(4)		# on appelle la fonction f avec comme argument 4 
							# et on récupère la valeur renvoyée dans la variable z
		\end{CodePythontexAlt}

\begin{remarque}[Variables locales et variables globales]
	
	\begin{itemize}
		\item les \textbf{variables locales} ne sont connues que dans la fonction
		\item les \textbf{variables globales} sont connues dans tout le programme
	\end{itemize}

\begin{CodePythontexAlt}[Centre,PremLigne=1, Lignes=false,]{}
print(x)   # ERREUR -> x est une variable locale (tout comme y et z)
\end{CodePythontexAlt}

\end{remarque}



\begin{easybox}{Vérification de la valeur d'une variable}{}

\begin{CodePythontexAlt}[Centre,PremLigne=1, Lignes=false,]{}
assert f(6) == 36     	# True -> aucune message 
assert f(6) == 12     	# False -> message d'erreur 
assert not(f(6) == 12)	# True -> aucun message

\end{CodePythontexAlt}
\end{easybox}

